% 分类目录：dl
\documentclass[12pt, a4paper]{article}
\usepackage[margin=2.5cm]{geometry}
\usepackage{ctex}
\usepackage{amsmath, amssymb}
\usepackage{listings}
\usepackage{xcolor}
\usepackage{tabularx}
\usepackage{hyperref}

\lstset{
    language=Python,
    basicstyle=\ttfamily\small,
    keywordstyle=\color{blue},
    commentstyle=\color{green!50!black},
    stringstyle=\color{red},
    showstringspaces=false,
    breaklines=true,
    numbers=left,
    numberstyle=\tiny\color{gray},
    frame=single
}

\title{知识点：正则化技术 (Regularization)}
\author{}
\date{}

\begin{document}
\maketitle

\section{核心定义}
\textbf{中文定义}：正则化（Regularization）是指通过修改学习算法，以降低测试误差（即泛化误差）为目的，而不仅仅是降低训练误差的方法。它主要通过向损失函数添加惩罚项或通过约束训练过程，来缓解深度学习模型中的过拟合（Overfitting）现象。

\textbf{英文定义}：Regularization refers to techniques used to modify a learning algorithm such that its test error (generalization error) is reduced, potentially at the cost of a slightly higher training error. It primarily mitigates overfitting in deep learning models by adding penalty terms to the loss function or constraining the training process.

\section{权重衰减 (Weight Decay / Norm Penalty)}
通过在原始损失函数 $L(\theta)$ 中加入参数范数惩罚项 $\Omega(\theta)$，来迫使参数变小：
$$ \tilde{L}(\theta) = L(\theta) + \lambda \Omega(\theta) $$

\subsection{L1 正则化 (Lasso)}
\textbf{惩罚项}：$\Omega(\theta) = \|\theta\|_1 = \sum |w_i|$。
\\ \textbf{特性}：产生稀疏解（Sparse Weights），某些不重要的权重会精确变为 0。常用于特征选择。

\subsection{L2 正则化 (Ridge)}
\textbf{惩罚项}：$\Omega(\theta) = \frac{1}{2} \|\theta\|_2^2 = \frac{1}{2} \sum w_i^2$。
\\ \textbf{特性}：迫使权重趋于较小的值（更平滑），但不会使其变为 0。在深度学习中也称为 \textbf{权重衰减}。

\section{暂退法 (Dropout)}
\textbf{原理}：在训练过程中，以概率 $p$ 随机将一部分神经元的激活值设为 0。
\\ \textbf{直观理解}：强制模型不能过度依赖某些特定的神经元路径，从而增强了鲁棒性。从统计学看，Dropout 可以被看作是一种廉价的集成学习（Ensemble Learning）方式，相当于对海量共享权重的子网络进行平均。
\\ \textbf{注意}：仅在训练时开启，预测时需关闭（并通常将权重乘以系数 $1-p$ 以保持期望一致）。

\section{提前停止 (Early Stopping)}
\textbf{策略}：在训练过程中监控验证集的误差，当验证集误差不再下降甚至开始上升时，提前终止训练。
\\ \textbf{优点}：不仅能防止过拟合，还减少了计算资源的消耗，是深度学习中最简单高效的正则化工具。

\section{数据增强 (Data Augmentation)}
通过旋转、翻转、缩放等手段人为增加训练样本的多样性。虽然它不直接修改网络结构，但它的泛化效果等同于强力的正则化。

\section{其他技术}
\begin{itemize}
    \item \textbf{标签平滑 (Label Smoothing)}：将 Hard Label（如 $[0, 1]$）转化为 Soft Label（如 $[0.05, 0.95]$），防止模型对预测值极度自信，有助于泛化。
    \item \textbf{批量归一化 (Batch Normalization)}：尽管其初衷是加速收敛，但其每个 Batch 的随机波动也带有微弱的正则化效果。
\end{itemize}

\section{关键代码片段 (PyTorch)}
\begin{lstlisting}
import torch.optim as optim
import torch.nn as nn

# 1. L2 正则化 (权重衰减)
# 在优化器中设置 weight_decay 参数
optimizer = optim.Adam(model.parameters(), lr=1e-3, weight_decay=1e-5)

# 2. Dropout 方法
# 直接在模型定义中加入 Dropout 层
class MyModel(nn.Module):
    def __init__(self):
        super().__init__()
        self.fc = nn.Linear(512, 256)
        self.dropout = nn.Dropout(p=0.5)
        
    def forward(self, x):
        return self.fc(self.dropout(x))
\end{lstlisting}

\section{深度面试题}

\textbf{问：为什么 L1 能产生稀疏权重，而 L2 只有缩小效果？}\\
答：从几何角度看，L1 的等值线是方形（在轴上有突出的尖角），梯度下降在更新时更容易在坐标轴上停留，导致分量为 0。而 L2 的等值线是圆形，权重只会无限接近 0 而不会真正达到 0。

\textbf{问：测试时如果忘了关 Dropout 会发生什么？}\\
答：模型的输出会具有随机性，即使对于相同的输入，每次预测结果也可能不同。此外，由于只有部分神经元工作，预测的量级通常会比预期小，导致预测准确率断崖式下跌。

\section{参考文献}
\begin{enumerate}
    \item Goodfellow, I., Bengio, Y., \& Courville, A. (2016). \textit{Deep Learning}. MIT Press (Chapter 7: Regularization for Deep Learning).
    \item Srivastava, N., et al. (2014). \textit{Dropout: A Simple Way to Prevent Neural Networks from Overfitting}. Journal of Machine Learning Research.
    \item Kukar, M., \& Kononenko, I. (2000). \textit{Cost-sensitive Learning with Label Smoothing}. AAAI.
\end{enumerate}

\end{document}
